\documentclass[12pt]{article}
\usepackage[utf8]{inputenc}
\usepackage[T1]{fontenc}
\usepackage{hyperref}
\title{Meroitic Language Resources for K-12 Curriculum Development}
\author{Compiled by Claude Sonnet 5.5 (claude-sonnet-5-5) for Andr\'as Kornai}
\date{}
\begin{document}
\maketitle

\section{Introduction}
Meroitic (ISO 639-3 xmr) was the language of the Kingdom of Kush in ancient Nubia. It is extinct and has no known descendants. The script was deciphered by F.~Ll.~Griffith in 1911, so the texts can be read aloud. The language itself is only partly understood: the syntax of nominal groups is fairly well known, but verbal morphology is still poorly understood. The main specialist in the field is Claude Rilly, who classifies Meroitic as Northern East Sudanic within Nilo-Saharan. No living speaker community exists, so a K-12 programme could only teach the script and the history, with a small amount of reconstructed vocabulary.

\section{Primary Educational Resources}
No school-level primers or textbooks for Meroitic were found.

\section{Linguistic and Reference Materials}
\begin{itemize}
\item Rilly and de Voogt, \emph{The Meroitic Language and Writing System} (Cambridge University Press): the main monograph. The search returned only the publisher's excerpt. \url{https://assets.cambridge.org/97811070/08663/excerpt/9781107008663_excerpt.htm}
\item The same excerpt as a PDF. \url{https://assets.cambridge.org/97811070/08663/excerpt/9781107008663_excerpt.pdf}
\item Claude Rilly, talk on Meroitic phonology, syntax and linguistic family (University of Chicago). \url{https://voices.uchicago.edu/lvc/2024/03/05/claude-rilly-cnrs-ephe-paris-sorbonne-meroitic-phonology-syntax-and-linguistic-family}
\item Article in Dotawo (a journal of Nubian studies) from the search results. \url{https://journalpub.escholarship.org/dotawo/article/53835/galley/40735/download}
\item Article in Mitteilungen der Sudanarch\"aologischen Gesellschaft. \url{https://journals.ub.uni-heidelberg.de/index.php/mittsag/article/view/90327/84978}
\end{itemize}

\section{Religious and Community Resources}
No living community or scripture resources exist. Sudan History has an overview of the Meroitic writing system. \url{https://sudanhistory.org/writing_02_meroitic.htm}

\section{Digital and Audio Resources}
A 2025 computational study treats Meroitic decipherment as a machine-learning task. \url{https://preview.aclanthology.org/moar-dois/2025.alp-1.11.pdf} A general-audience article on undeciphered ancient languages is at \url{https://amp.livescience.com/59851-ancient-languages-not-yet-deciphered.html}. No audio or software suitable for classroom use was found.

\section{Curriculum Development Context}
Very little suitable material exists. A realistic school unit would cover the history of Kush, the syllabic script with its default vowel /a/, and the way decipherment works. It would not offer conversational use of the language.
\end{document}